\label{ch:appendix-additional-results}

This appendix collects the supplementary breakdowns behind
Section~\ref{sec:results-cost}, which reports what iterative refinement
costs and where that cost comes from. The splits below are descriptive
detail rather than a research question in their own right: none of them
are used to rank the three models beyond what
Table~\ref{tab:setup-model-cost} already establishes, only to show where
the money actually went. Figures are computed directly from each task's
\texttt{llm\_cost\_ledger.json} (its per-iteration and per-call-type cost
breakdown) via \texttt{scripts/analysis/cost\_breakdown.py}, cross-checked
against the per-task totals in
\texttt{results\_aggregate\_lan/cells.csv}. Both are read from
\texttt{results\_lan}, the network-consistent tree every reported number
in Chapter~\ref{ch:results} comes from
(Section~\ref{sec:results-load-profile-validation}).

\paragraph{A note on measurement.} Spend is what the providers actually
billed, so it includes the overhead of running the campaign rather than an
idealised price. Nine times across seven tasks a run was stopped and later
resumed; because the conversation is replayed from disk while the
provider-side prompt cache has gone cold, the accumulated history is
re-billed at the uncached rate on the first calls after each resume. That
overcharge is \$7.99, or 2.6\% of total spend, and it falls between 1.6\%
and 2.8\% for every model, so it inflates the absolute figures slightly and
leaves the relative comparison intact. The resumes affect nothing else:
the replayed conversation is identical, and all 63 tasks ran their full
eleven iterations.

\section{LLM Cost Breakdown}
\label{sec:appendix-cost}

\paragraph{Spend by framework and scenario.} Tables~\ref{tab:appendix-cost-env}
and~\ref{tab:appendix-cost-scenario} give the per-task means behind the
brief treatment in Section~\ref{sec:results-cost}. Neither split produces a
consistent ordering: the cheapest framework differs for all three models,
and the scenario spread is flat apart from \textsc{Petstore}.

\begin{table}[htbp]
  \centering
  \caption{Mean LLM spend (USD) per task, by framework and model. Each
    cell pools that framework's seven tasks for that model.}
  \label{tab:appendix-cost-env}
  \small
  \begin{tabular}{@{}lrrrr@{}}
    \toprule
    \textbf{Framework} & \textbf{\texttt{claude-opus-4-8}} & \textbf{\texttt{gpt-5.5}} & \textbf{\texttt{glm-5.2}} & \textbf{All} \\
    \midrule
    \textsc{Go-net-http}  & \$6.48 & \$7.14 & \$1.77 & \$5.13 \\
    \textsc{Python-Flask} & \$5.68 & \$8.04 & \$1.21 & \$4.98 \\
    \textsc{Rust-Actix}   & \$6.93 & \$6.06 & \$1.09 & \$4.69 \\
    \bottomrule
  \end{tabular}
\end{table}

\begin{table}[htbp]
  \centering
  \caption{Mean LLM spend (USD) per task, by scenario and model, ordered
    by the pooled mean. Each cell pools that scenario's three tasks for
    that model.}
  \label{tab:appendix-cost-scenario}
  \small
  \resizebox{\linewidth}{!}{%
  \begin{tabular}{@{}lrrrr@{}}
    \toprule
    \textbf{Scenario} & \textbf{\texttt{claude-opus-4-8}} & \textbf{\texttt{gpt-5.5}} & \textbf{\texttt{glm-5.2}} & \textbf{All} \\
    \midrule
    \textsc{ParcelPinLockerPickup}         & \$5.19 & \$5.44  & \$1.34 & \$3.99 \\
    \textsc{ClickCount}                    & \$6.95 & \$5.40  & \$1.00 & \$4.45 \\
    \textsc{TransitPulseDelayReporter}     & \$5.72 & \$6.17  & \$1.62 & \$4.50 \\
    \textsc{BranchWeave}                   & \$6.59 & \$5.86  & \$1.48 & \$4.64 \\
    \textsc{Recipes}                       & \$6.68 & \$6.75  & \$0.92 & \$4.78 \\
    \textsc{SplitNestSharedExpenseLedger}  & \$6.14 & \$6.68  & \$1.94 & \$4.92 \\
    \textsc{Petstore}                      & \$7.28 & \$13.26 & \$1.20 & \$7.25 \\
    \bottomrule
  \end{tabular}%
  }
\end{table}

\paragraph{Cost to reach a working baseline.} Iteration~000's cost is the
first LLM call of a task: a single baseline code-and-spec generation
attempt, before any runtime feedback exists to refine against. For tasks
that start with a healthy baseline (Section~\ref{sec:results-rq1}), that
one call is the entire cost of getting to a working deployment.
For tasks that start dead, reaching a working deployment instead costs
whatever iterations~000 through the first iteration with positive
goodput together spend, since a dead baseline is only useful once
refinement has recovered it. Table~\ref{tab:appendix-cost-baseline}
reports both, per model.

\begin{table}[htbp]
  \centering
  \caption{Mean LLM cost (USD) to reach a working baseline, split by
    whether the task started healthy (cost of \texttt{iteration-000}
    alone) or started dead and had to recover (cumulative cost through
    the first iteration with positive goodput). $n$ is the number of
    tasks in each split; the two do not always sum to 21 because two
    tasks (Section~\ref{sec:results-rq1}) never recover a positive
    baseline and are excluded from the dead-baseline column.}
  \label{tab:appendix-cost-baseline}
  \small
  \begin{tabular}{@{}lrrrr@{}}
    \toprule
    & \multicolumn{2}{c}{\textbf{Started healthy}} & \multicolumn{2}{c}{\textbf{Started dead, recovered}} \\
    \textbf{Model} & \textbf{Mean \$} & \textbf{$n$} & \textbf{Mean \$} & \textbf{$n$} \\
    \midrule
    \texttt{claude-opus-4-8} & \$0.74 & 12 & \$2.42 & 8 \\
    \texttt{gpt-5.5}         & \$1.03 & 14 & \$4.54 & 7 \\
    \texttt{glm-5.2}         & \$0.11 & 7  & \$0.39 & 13 \\
    \bottomrule
  \end{tabular}
\end{table}

Recovering a dead baseline costs more than an already-healthy one for
every model, three to five times as much on average, which is the
expected shape: recovery is itself several refinement iterations rather
than the single baseline-generation call a healthy task needed. The gap
is proportionally similar across models despite their very different
absolute price points, so this cost is a consequence of how many extra
iterations recovery took, not of one model being disproportionately
expensive to recover with.

\paragraph{Cost by refinement lever.} Every refinement iteration spends
on a refinement-decision call (choosing between the code and deployment
lever, Section~\ref{sec:methods-workflow}) and then on whichever lever it
picked. Table~\ref{tab:appendix-cost-lever} reports the mean cost of each,
per model, computed as total spend on that call type divided by the
number of steps of that kind actually taken (Table~\ref{tab:results-rq3-funnel}
gives those step counts).

\begin{table}[htbp]
  \centering
  \caption{Mean LLM cost (USD) per refinement-decision call, per
    code-refinement step, and per deployment-refinement step, by model.}
  \label{tab:appendix-cost-lever}
  \small
  \resizebox{\linewidth}{!}{%
  \begin{tabular}{@{}lrrrr@{}}
    \toprule
    \textbf{Model} & \textbf{\$ / decision} & \textbf{\$ / code step} & \textbf{\$ / deployment step} & \textbf{Ratio} \\
    \midrule
    \texttt{claude-opus-4-8} & \$0.24 & \$0.58 & \$0.22 & $2.6\times$ \\
    \texttt{gpt-5.5}         & \$0.18 & \$0.68 & \$0.25 & $2.7\times$ \\
    \texttt{glm-5.2}         & \$0.07 & \$0.11 & \$0.05 & $2.1\times$ \\
    \bottomrule
  \end{tabular}%
  }
\end{table}

A code-refinement step costs two to three times as much as a
deployment-refinement step for every model, which tracks what the two calls
actually do: code refinement rewrites and resubmits application source,
while deployment refinement adjusts a small, structured \texttt{spec.yaml}
(Section~\ref{sec:methods-spec}). Section~\ref{sec:results-rq3} reports
that code refinement also fails outright more often than deployment refinement
before ever producing a usable comparison; on top of that risk, it is
consistently the more expensive lever to attempt, for all three models.

\paragraph{How spend accumulates over the budget.}
Table~\ref{tab:appendix-cost-iteration} gives the mean per-task spend at
each iteration index. The baseline call is the single most expensive
iteration for every model, but only 13\% of a task's total; the ten
refinement iterations that follow are close to flat, so total spend is
close to linear in the iteration budget. This is what makes the budget,
rather than any per-call property, the parameter that sets the price of a
run.

\begin{table}[htbp]
  \centering
  \caption{Mean LLM spend (USD) per task at each iteration index, and the
    cumulative share of a task's total spend reached by that index
    (pooled over all three models).}
  \label{tab:appendix-cost-iteration}
  \small
  \resizebox{\linewidth}{!}{%
  \begin{tabular}{@{}lrrrrr@{}}
    \toprule
    \textbf{Iteration} & \textbf{\texttt{claude-opus-4-8}} & \textbf{\texttt{gpt-5.5}} & \textbf{\texttt{glm-5.2}} & \textbf{All} & \textbf{Cum.\ share} \\
    \midrule
    000 & \$0.64 & \$1.19 & \$0.12 & \$0.65 & 13\% \\
    001 & \$0.40 & \$0.45 & \$0.12 & \$0.32 & 20\% \\
    002 & \$0.48 & \$0.54 & \$0.10 & \$0.37 & 27\% \\
    003 & \$0.59 & \$0.51 & \$0.08 & \$0.39 & 35\% \\
    004 & \$0.64 & \$0.49 & \$0.07 & \$0.40 & 43\% \\
    005 & \$0.63 & \$0.49 & \$0.11 & \$0.41 & 52\% \\
    006 & \$0.58 & \$0.78 & \$0.13 & \$0.49 & 62\% \\
    007 & \$0.64 & \$0.49 & \$0.12 & \$0.42 & 70\% \\
    008 & \$0.60 & \$0.62 & \$0.18 & \$0.47 & 80\% \\
    009 & \$0.61 & \$0.66 & \$0.15 & \$0.47 & 89\% \\
    010 & \$0.56 & \$0.87 & \$0.17 & \$0.53 & 100\% \\
    \bottomrule
  \end{tabular}%
  }
\end{table}
